\documentclass[../infufrgs/ppgc-tese.tex]{subfiles}
\begin{document}

This chapter reports the completed G3 study of the thesis: an empirical
characterization of what an activation-space adversarial detector, built on the
sparse-feature pipeline of Chapter~\ref{chpt:circuitry}, actually measures in a
Vision Transformer classifier. The pre-registered protocol, hypotheses, and gates
were declared in Chapter~\ref{chpt:circuitry}; here we execute them and read the
result. That result is deflationary, and deliberately so: to the resolution of a
magnitude-sensitive linear probe, the detector reads off-manifold displacement
\emph{magnitude}, not adversarial-specific semantic structure, and the sparse
autoencoder (SAE) is not load-bearing for the detection. We summarize this as
\emph{magnitude, not meaning}.

\section{The confound: adversarial versus perturbed}
\label{sec:adv-confound}
A recurring hope in interpretability is that a model's own internal
representations betray an attack --- that somewhere in activation space there is a
signal for ``this input is adversarial'' --- and that SAEs, by decomposing
activations into interpretable features, can expose it. On this reading,
adversarial detection becomes a form of introspection: the model, read through its
features, ``knows'' it is under attack.

This hope rests on a confound that is rarely isolated. Every adversarial example is
two things at once: it is \emph{adversarial} (crafted to change the output) and it
is a \emph{perturbation} (the activation has moved off its clean location). A
detector that fires on adversarial inputs may be reading either property --- the
adversarial structure, or merely the displacement. The two are entangled in any
standard evaluation, because the only negatives are clean inputs, which differ from
adversarial ones in both respects at once. Disentangling them requires a control
that standard detection setups omit: a benign perturbation of the \emph{same
magnitude} as the attack.

Our question is therefore one of characterization, not deployment. We do not ask
whether an activation-space detector can be built --- it plainly can --- but what
such a detector measures. We study a ViT classifier (ViT-B/16, with ViT-L/16 as a
scale check), not a vision--language model, and probe its MLP activations with a
linear detector under the leakage-free, grouped cross-validation protocol of
Chapter~\ref{chpt:circuitry}.

\section{Instrument: the residual decomposition and its controls}
\label{sec:adv-setup}
The full protocol --- models, pre-registered hypotheses, phase plan, and design
decisions --- is specified in Chapter~\ref{chpt:circuitry}; we recall here only what
is needed to read the results. Every activation is the MLP output of a transformer
block, mean-pooled over the $197$ tokens, at blocks $\{0,3,6,9,10,11\}$ of ViT-B/16.
We decompose each activation into an SAE reconstruction and a residual,
\begin{equation}
  x \;\approx\; \hat{x} + r,
  \qquad \hat{x} = \mathrm{dec}(z), \quad z = \mathrm{ReLU}(\mathrm{enc}(x)),
  \qquad r = x - \hat{x},
  \label{eq:adv-decomp}
\end{equation}
where $x$ is the raw activation, $z$ the sparse latents, $\hat{x}$ the
reconstruction, and $r$ the residual. Critically, the SAE is trained on \emph{clean
activations only}, so $\hat{x}$ is a projection onto the clean-data manifold and $r$
is, by construction, the off-manifold component the clean dictionary cannot explain.
Across the six layers the SAE explains between $0.752$ and $0.948$ of activation
variance.

The logic of the decomposition is a competition. We evaluate detection AUC
separately for the four vector representations $x$, $z$, $\hat{x}$, $r$ and for the
scalar residual norm $\lVert r\rVert_2$. If interpretable SAE features carried the
adversarial signal, the latents $z$ should win; if instead the signal is mere
off-manifold displacement, the scalar $\lVert r\rVert_2$ --- a pure magnitude ---
should already capture most of it. The residual-norm row is thus the decisive
control.

We contrast two attacks that differ chiefly in magnitude. Our high-$L_2$ attack is
untargeted FGSM with an $L_\infty$ budget \cite{goodfellow2015fgsm}, canonical
setting $\varepsilon{=}0.012$ (swept over
$\varepsilon\in\{0.001,\dots,0.032\}$ for the magnitude test). Our low-$L_2$ anchor
is a steganographic text watermark at opacity $0.03$, which carries a near-zero
$L_2$ payload. To separate ``adversarial'' from merely ``perturbed'', we add four
benign perturbation families calibrated to the same $L_2$ as FGSM at
$\varepsilon{=}0.012$: Gaussian noise, Gaussian blur, JPEG compression, and random
pixel dropout. Every detection AUC is computed by grouped $5$-fold
cross-validation keyed to the \emph{source image} --- all perturbations of an image
share a fold, so the detector is never tested on a near-duplicate of a training
example --- with standardized features and a converged logistic regression; each
binary task has $n{=}3200$ examples ($1600$ clean, $1600$ perturbed).

\section{What the detector reads: the decomposition}
\label{sec:adv-decomp-results}
Table~\ref{tab:adv-decomp} reports the decomposition AUCs and reads as a four-link
chain.

\begin{table}[ht]
  \centering
  \caption{Detection AUC (grouped 5-fold CV, $n{=}3200$) by representation and
  layer. Raw activation $x$ saturates for FGSM at \emph{every} layer; the SAE
  latents $z$ never improve on $x$; steganography is near chance throughout.}
  \label{tab:adv-decomp}
  \begin{tabular}{llcccc}
    \toprule
    Layer & Task & $x$ & $z$ & $\hat{x}$ & $r$ \\
    \midrule
    \multirow{2}{*}{mlp\_0}  & FGSM  & 1.000 & 0.996 & 1.000 & 1.000 \\
                             & stego & 0.675 & 0.539 & 0.541 & 0.677 \\
    \multirow{2}{*}{mlp\_3}  & FGSM  & 1.000 & 1.000 & 1.000 & 1.000 \\
                             & stego & 0.685 & 0.604 & 0.604 & 0.664 \\
    \multirow{2}{*}{mlp\_6}  & FGSM  & 1.000 & 1.000 & 1.000 & 1.000 \\
                             & stego & 0.549 & 0.531 & 0.537 & 0.542 \\
    \multirow{2}{*}{mlp\_9}  & FGSM  & 1.000 & 0.999 & 0.998 & 1.000 \\
                             & stego & 0.530 & 0.513 & 0.524 & 0.530 \\
    \multirow{2}{*}{mlp\_10} & FGSM  & 1.000 & 1.000 & 1.000 & 1.000 \\
                             & stego & 0.543 & 0.519 & 0.532 & 0.540 \\
    \multirow{2}{*}{mlp\_11} & FGSM  & 1.000 & 0.997 & 1.000 & 1.000 \\
                             & stego & 0.544 & 0.513 & 0.524 & 0.542 \\
    \bottomrule
  \end{tabular}
\end{table}

\emph{First, raw activations already detect FGSM perfectly.} The raw representation
$x$ reaches AUC $1.000$ at \emph{every} layer, including the first MLP block
(\texttt{mlp\_0}).

\emph{Second, the SAE is not load-bearing.} The latent representation $z$ never
exceeds $x$, and is strictly below it at four of six layers ($z{=}0.996$ at
\texttt{mlp\_0}, $0.999$ at \texttt{mlp\_9}, $0.997$ at \texttt{mlp\_11}); the
reconstruction $\hat{x}$ behaves the same. Because the SAE is a lossy bottleneck it
can only discard detection information, never manufacture it --- so the interpretable
features carry no detection signal that the raw activation lacks.

\emph{Third, the signal is low-level, not late-semantic.} A perfect AUC at
\texttt{mlp\_0}, before nearly any transformer computation, refutes the hypothesis
that the detector reads a late ``I am under attack'' representation. This layer-0
outcome was pre-registered as an anti-condition in Chapter~\ref{chpt:circuitry}.

\emph{Fourth, the low-$L_2$ attack is invisible.} Steganography separates from clean
only weakly at the shallowest layer ($x{=}0.675$ at \texttt{mlp\_0}) and decays
toward chance with depth ($\approx0.54$ by \texttt{mlp\_11}); on every SAE
representation it sits at chance throughout. Steganography is the low-$L_2$ anchor of
the magnitude argument, not an incidental weak result.

The scalar residual norm $\lVert r\rVert_2$ --- the pure-magnitude control ---
sharpens the reading. For FGSM it alone reaches AUC $0.984$, $1.000$, $0.929$ at
\texttt{mlp\_0/3/6}, essentially matching the full vector residual $r$ ($=1.000$),
but falls to $0.660$ and $0.586$ at \texttt{mlp\_9} and \texttt{mlp\_10} while $r$
stays at $1.000$. So at shallow layers detection is well described as pure
off-manifold displacement magnitude, whereas at deep layers the displacement
\emph{direction} carries information beyond its scalar norm --- though still not
through interpretable SAE features, since $z$ never beats $x$.

\section{Magnitude drives detection: H1 and H2}
\label{sec:adv-h1h2}

\subsection{H1 --- detectability scales with perturbation magnitude}
Table~\ref{tab:adv-h1} shows detection AUC rising monotonically with the
perturbation budget, from $0.670$ at $\varepsilon{=}0.001$ to $1.000$, saturating
for $\varepsilon\geq0.012$. This dependence is the signature of a magnitude detector,
and it fails the pre-registered H1 invariance gate on all three of its criteria: the
gate requires a near-zero $\varepsilon$--AUC correlation ($|r|<0.3$), low AUC
variance ($<10^{-3}$), and mean AUC $>0.95$ over $\varepsilon\leq0.012$; the sweep
violates each, its mean over the five budgets $\varepsilon\leq0.012$ being
$\approx0.88$, below the floor.

\begin{table}[ht]
  \centering
  \caption{H1: FGSM detection AUC versus perturbation budget $\varepsilon$ under the
  corrected protocol. AUC rises with $\varepsilon$ and saturates by
  $\varepsilon\geq0.012$; the pre-registered H1 invariance gate fails on all three
  criteria.}
  \label{tab:adv-h1}
  \begin{tabular}{lcccccccc}
    \toprule
    $\varepsilon$ & 0.001 & 0.002 & 0.004 & 0.008 & 0.012 & 0.016 & 0.024 & 0.032 \\
    \midrule
    AUC & 0.670 & 0.785 & 0.932 & 0.996 & 1.000 & 1.000 & 1.000 & 1.000 \\
    \bottomrule
  \end{tabular}
\end{table}

This failure must be read in the right direction. Saturation for
$\varepsilon\geq0.012$ \emph{removes} discriminative power from the gate rather than
adding it: once AUC is pinned at the ceiling, the gate can no longer separate a
magnitude detector from a structure detector in that regime. We therefore do not
claim the corrected protocol ``strengthens'' the negative result; it locates the
informative, sub-saturation regime ($\varepsilon\approx0.004$, AUC $0.932$) that a
clinching test should use.

\subsection{H2 --- the signal is not adversarial-specific}
Table~\ref{tab:adv-h2} pits FGSM against four benign perturbation families of
comparable magnitude. FGSM is detected at AUC $0.9999$ --- but so are the benign
families: JPEG at $0.939$, random dropout $0.921$, Gaussian blur $0.910$, and
Gaussian noise $0.847$. The FGSM--benign gap is only $0.0605$, far short of the
pre-registered $0.3$ threshold, and the strongest benign family ($0.939$) far
exceeds the $0.7$ ceiling; the H2 specificity gate fails on both criteria.

\begin{table}[ht]
  \centering
  \caption{H2: detection AUC for FGSM versus matched-magnitude benign perturbations.
  The FGSM--benign gap is small ($\approx0.06$) and the specificity gate fails.}
  \label{tab:adv-h2}
  \begin{tabular}{lc}
    \toprule
    Family & AUC \\
    \midrule
    FGSM (adversarial)       & 0.9999 \\
    JPEG (benign)            & 0.939 \\
    random\_dropout (benign) & 0.921 \\
    gaussian\_blur (benign)  & 0.910 \\
    gaussian\_noise (benign) & 0.847 \\
    \midrule
    FGSM\,$-$\,max(benign) gap & 0.0605 \\
    \bottomrule
  \end{tabular}
\end{table}

The reading is that the detector is not reading ``adversarial'' but ``displaced'': a
benign perturbation of similar magnitude is nearly as detectable. Two flanks of this
run still restrict the strongest form of the claim. First, this run does not emit the
realized per-family $L_2$, so ``matched-$L_2$'' is here the design intent written
into the perturbation configs rather than a verified measurement
\pendingvtwo{the v2 re-run emits per-family $L_2$ within $\pm20\%$, verifying the
match}. Second, FGSM is saturated at $\approx1.0$, so the gap is measured against a
ceiling rather than in the informative sub-saturation regime
\pendingvtwo{v2 re-runs H2 at $\varepsilon\approx0.004$ so the FGSM--benign gap has
headroom}. Both flanks are closed by the planned v2 re-run; the direction of the
present evidence already supports non-specificity, and the clinching, fully
magnitude-controlled version follows once v2 lands.

\section{Analysis: the specificity ladder and a reconciliation with SAEgis}
\label{sec:adv-analysis}
It helps to organize ``how specific is the detected signal?'' as a ladder.
\emph{Degree~1} is specificity to the \emph{attack method} (features for one attack
family differ from another). \emph{Degree~2} is specificity to the \emph{adversarial
class} (adversarial features differ from a benign perturbation of equal magnitude).
\emph{Degree~3} is \emph{generic displacement} (the detector fires for anything
off-manifold, attack or not). A matched-$L_2$ benign control is precisely the
instrument that separates Degree~2 from Degree~3; without it, an evaluation cannot
distinguish an adversarial-specific detector from a displacement detector.

This framing positions our study relative to SAEgis \cite{wang2026saegis}, a
contemporaneous SAE ``firewall'' for adversarial detection in vision--language
models. SAEgis is a \emph{deployment} paper --- can a firewall be built and made to
generalize? --- whereas ours is a \emph{characterization} paper --- what does the
detector read? We contradict none of their measurements; where the two overlap, their
results corroborate ours. Their dense (no-SAE) baseline ties their SAE in-domain,
independent support that the SAE is not load-bearing for in-distribution detection of
the perturbation; their reconstruction-error anomaly score shows a negligible
clean-versus-adversarial difference, consistent with our finding that the scalar
residual norm is a weak signal relative to the residual \emph{vector}; and their
earliest-layer detector is their most robust across attacks, corroborating that the
cue is low-level rather than late-semantic. SAEgis's own cross-attack analysis
already descends from Degree~1; the matched-$L_2$ benign control we add is the
experiment that extends the descent toward Degree~3. Their cross-domain
generalization benefit is real and orthogonal; we do not claim SAEs are useless, only
that the SAE adds nothing to in-distribution detection of the perturbation, and that
the detection is magnitude-driven. The contradiction we do raise is aimed narrowly at
the \emph{strong} reading of ``attack-relevant'' (semantic specificity to the
adversarial class); the \emph{weak} reading (useful for detection) stands.

\section{Methodological case studies}
\label{sec:adv-bugs}
Two implementation bugs, caught and fixed during this work, are instructive enough to
report: each would have produced a confident but false positive, and each maps to a
cheap sanity check that activation-based detection studies should adopt. We present
them as such, not as confessions.

\paragraph{The clip bug: an out-of-distribution artifact masquerading as
magnitude-invariant detection.} Our FGSM routine originally applied
$\mathrm{clamp}(x+\varepsilon\,\mathrm{sign}(\nabla),0,1)$ directly to
ImageNet-\emph{normalized} images, whose values span roughly $[-2.2,2.6]$. Clipping
to $[0,1]$ in that space squashes most pixels, so the ``attack'' was dominated by
clip distortion rather than the $\varepsilon$ perturbation: at $\varepsilon{=}0.012$
the measured per-image $L_2$ was $\approx284$, versus the $\approx20$ expected from
the $\varepsilon$ arithmetic. The clipped image left the normalized manifold entirely
--- a synthetic, $\varepsilon$-independent OOD signal that the detector simply read
off, producing a \emph{flat} pre-fix $\varepsilon$-sweep (a false confirmation of a
magnitude-invariant detector). The fix clips where it is physically meaningful ---
denormalize, clip in $[0,1]$, renormalize --- after which the per-image $L_2$ at
$\varepsilon{=}0.012$ is $20.51$, consistent with theory, and the sweep recovers the
monotone curve of Table~\ref{tab:adv-h1}. The transferable lesson: measure the
perturbation $L_2$ and compare it to the $\varepsilon$ arithmetic \emph{before}
reporting any detection AUC.

\paragraph{The first-key bug: a layer-provenance error in SAE training.} The SAE
trainer selected its target activations as the first non-label key of the saved
activation file. With forward hooks registered in the order
$[\texttt{layer\_10},\texttt{layer\_11}]$, that first key was always
\texttt{layer\_10}, while the \texttt{target\_layer} configuration field was dead
code with no effect. The project's ``canonical'' detection number was thus computed
on \texttt{mlp\_10} alone, not the \texttt{mlp\_11} the documentation claimed --- a
provenance error, not an invalidation of the underlying measurement. We fixed it by
plumbing an explicit target-layer argument through SAE training and latent encoding,
which also removes the silent dependence on hook registration order. The transferable
lesson: never let an analysis target be selected implicitly by container ordering;
make it an explicit, logged argument.

\section{Limitations}
\label{sec:adv-limitations}
Our claims are bounded by our instrument, and we state the boundaries plainly.

\paragraph{A null is a null of this probe.} Detection here is a \emph{linear}
classifier on \emph{mean-pooled} activations. Both choices limit what a null can
mean. On the same SAE latents at \texttt{mlp\_11}, replacing the linear detector with
a small MLP raised AUC from $0.547$ to $0.685$ --- roughly $0.14$ of AUC that the
linear probe could not see --- and mean-pooling over the $197$ tokens erases any
token-localized structure. We therefore cannot claim that fine-grained adversarial
structure is \emph{absent}; only that it is invisible to a magnitude-sensitive
linear, mean-pooled probe under matched conditions. A non-linear, non-pooled probe
would be needed to move ``we do not see it'' toward ``it is not there.''

\paragraph{Saturation weakens the gates.} FGSM detection is pinned at AUC
$\approx1.0$ for $\varepsilon\geq0.012$. This ceiling removes discriminative power
from both pre-registered gates: H1 cannot distinguish magnitude- from
structure-dependence once saturated, and the H2 FGSM--benign gap is measured against
the ceiling rather than in the informative sub-saturation regime.

\paragraph{Matched-$L_2$ measurement and attack coverage.} The strong,
adversarial-specificity-controlled form of the claim rests on emitting the realized
per-family $L_2$ and measuring the FGSM--benign gap at a non-saturated
$\varepsilon\approx0.004$ \pendingvtwo{the v2 re-run supplies both}. Beyond FGSM and
steganography, PGD \cite{madry2018pgd} (pre-registered as the most rigorous H1 test)
and a Mahalanobis-on-raw-activations baseline \cite{lee2018mahalanobis} are planned
but not yet run. The H1/H2 results rest on a single seed, the single canonical layer
(\texttt{mlp\_10}), and one SAE architecture.

\paragraph{Scope of the system.} We study a ViT classifier and make no claim about
vision--language models, which we do not test.

\section{Summary}
\label{sec:adv-summary}
Under an SAE residual decomposition, raw activations detect a high-$L_2$ attack
perfectly at every layer, the SAE latents add nothing, the signal is already present
at layer~0, and a near-zero-$L_2$ attack is undetectable; benign perturbations of
comparable magnitude are nearly as detectable as the attack. To the resolution of a
magnitude-sensitive linear probe, activation-space adversarial detection in this ViT
classifier is a displacement-magnitude detector, and the interpretable features are
not load-bearing for it. This is the sense in which detection here is
\emph{magnitude, not meaning}. The deflation of ``meaning'' --- the SAE is not
load-bearing, and the cue is low-level rather than late-semantic --- is established
now; the fully magnitude-controlled form of ``magnitude'' is clinched by the v2
re-run that verifies the $L_2$ match and measures the gap at a non-saturated budget.
This characterization directly informs the thesis-level claim on shift/attack
diagnostics (G3, Chapter~\ref{chpt:circuitry}): internal sparse-feature signatures
match, rather than surpass, dense-activation baselines, because both read the same
low-level displacement.

\end{document}
